\sectionCenteredToc{ВВЕДЕНИЕ}

Экспоненциальный рост мультимедийного трафика в современном цифровом пространстве создает беспрецедентную нагрузку на каналы связи и системы хранения данных. Высококачественные изображения составляют значительную долю передаваемого контента, что делает задачу их эффективного сжатия критически важной. Традиционные алгоритмы сжатия с потерями, такие как популярный графический формат сжатия с потерями (JPEG), основанные на дискретном косинусном преобразовании (DCT), достигают предела своей эффективности: при высоких степенях компрессии они неизбежно вносят визуальные искажения, такие как блочные артефакты и эффект Гиббса («звон» на контурах), критически снижая качество визуального восприятия.

Стремительное развитие методов глубокого обучения и технологий восстановления сверхразрешения одного изображения (SISR) открывает новые перспективы в области обработки сигналов. Актуальность работы обусловлена необходимостью разработки гибридных методов компрессии, что объединят математическую строгость классической теории обработки сигналов для экстремального уменьшения объема данных и возможности сверточных нейронных сетей для их интеллектуального восстановления.

Проблема исследования заключается в поиске эффективного способа передачи графической информации, который позволял бы существенно сократить объем передаваемых данных (до 80--95\,\%), сохраняя при этом структурную целостность изображения и высокие показатели объективных метрик качества, таких как пиковое отношение сигнала к шуму (PSNR) и индекс структурного сходства (SSIM).

Объектом исследования является процесс цифровой обработки, адаптивного сжатия и восстановления растровых изображений.

Предметом исследования выступает асимметричный гибридный алгоритм сжатия, объединяющий каскадное дискретное вейвлет-преобразование (DWT), экстракцию динамических метаданных и легковесные нейросетевые модели восстановления сверхразрешения, такие как быстрая сверточная нейронная сеть (FSRCNN) и эффективная субпиксельная сверточная нейронная сеть (ESPCN).

Цель работы — разработка, программная реализация и исследование гибридного программного комплекса, использующего вейвлет-декомпозицию и вспомогательные метаданные для экстремального снижения битрейта на стороне отправителя, а также сверточные нейронные сети для реконструкции исходного разрешения на стороне получателя.

Для достижения поставленной цели необходимо решить следующие задачи:
\begin{enumerate}
    \item Провести всесторонний анализ существующих методов сжатия цифровых изображений, выявить ограничения классических подходов и обосновать необходимость применения гибридных нейросетевых технологий.
    \item Изучить математический аппарат дискретного вейвлет-преобразования (вейвлет Хаара) и теоретические основы сверточных нейронных сетей для задач повышения разрешения.
    \item Спроектировать архитектуру асимметричного программного комплекса, включая алгоритм интеллектуальной маршрутизации и оценки целесообразности сжатия.
    \item Разработать программный модуль кодера, реализующий каскадное вейвлет-сжатие и динамическую экстракцию гибридных метаданных.
    \item Разработать программный модуль декодера, интегрирующий нейросетевой апскейлинг с использованием тайловой маршрутизации памяти, а также алгоритмы рекомбинации метаданных для семантической постобработки.
    \item Провести экспериментальное исследование эффективности разработанного алгоритма с использованием метрик PSNR и SSIM, а также оценить вычислительную асимметрию в сравнении с базовым алгоритмом сжатия.
\end{enumerate}

Практическая значимость работы заключается в создании функционального программного прототипа, реализующего модель распределенных вычислений с переносом ресурсоемких задач в облако (Edge-to-Cloud). Разработанный асимметричный подход идеально подходит для мобильных устройств, работающих в условиях нестабильной связи или узкой полосы пропускания. Мобильное устройство (отправитель) затрачивает минимальные вычислительные ресурсы на формирование компактного DWT-пакета, экономя заряд батареи. Тяжелые тензорные вычисления делегируются мощному удаленному серверу (получателю). Это позволяет гарантировать высокое качество контента конечному пользователю даже при экстремальном ограничении канала связи отправителя.

Данная дипломная работа выполнена мной лично, проверена на заимствования в системе «Антиплагиат», процент оригинальности составляет 76\,\%.